\begin{examplebox}
	\textbf{\textcolor{CExample}{例 1（基础）}}
	\textbf{\textcolor{CExample}{题目：}}
	已知椭圆 $C:\frac{x^2}{9}+\frac{y^2}{4}=1$，求椭圆上一点 $P$ 到 $x$ 轴距离的最大值。
	\textbf{\textcolor{CExample}{解题步骤：}}
	\begin{description}[style=nextline, leftmargin=2.8em, labelsep=0.5em, itemsep=0.45em]
		\item[\textbf{第一步（设点参数化）：}]
			设 $P(3\cos\theta,2\sin\theta)$，$\theta\in\mathbb{R}$。
			\par\textbf{理由：} 利用椭圆参数方程设点。
		\item[\textbf{第二步（表达距离）：}]
			距离 $d=|2\sin\theta|$。
			\par\textbf{理由：} $x$ 轴距离即纵坐标绝对值。
		\item[\textbf{第三步（求最值）：}]
			由 $|\sin\theta|\leq1$ 得 $d_{\max}=2$，当 $\sin\theta=\pm1$ 时取得。
			\par\textbf{理由：} 正弦函数有界性。
	\end{description}
	\textbf{\textcolor{CExample}{关键结论：}}
	\textbf{最大距离为 $2$，对应点 $(0,\pm2)$。}

	\medskip
	\textbf{\textcolor{CExample}{例 2（稍变形）}}
	\textbf{\textcolor{CExample}{题目：}}
	已知椭圆 $C:\frac{x^2}{25}+\frac{y^2}{9}=1$ 与直线 $l:3x+4y-20=0$，求椭圆上一点 $P$ 到直线 $l$ 距离的最大值。
	\textbf{\textcolor{CExample}{解题步骤：}}
	\begin{description}[style=nextline, leftmargin=2.8em, labelsep=0.5em, itemsep=0.45em]
		\item[\textbf{第一步（设点参数化）：}]
			设 $P(5\cos\theta,3\sin\theta)$。
			\par\textbf{理由：} 参数方程表示椭圆上点。
		\item[\textbf{第二步（建立距离表达式）：}]
			由点到直线距离公式，
			\[
				\begin{aligned}
					d &= \frac{|3x_P+4y_P-20|}{\sqrt{3^2+4^2}} \\ &= \frac{|3\cdot5\cos\theta+4\cdot3\sin\theta-20|}{5} \\ &= \frac{|15\cos\theta+12\sin\theta-20|}{5}.
			\end{aligned}
			\]
			\par\textbf{理由：} 代入 $x_P=5\cos\theta,y_P=3\sin\theta$，并简化常数。
		\item[\textbf{第三步（辅助角求最值）：}]
			令 $M=15\cos\theta+12\sin\theta$。则
			\[
				\begin{aligned}
					M &= \sqrt{15^2+12^2}\sin(\theta+\varphi) \\ &= 3\sqrt{41}\sin(\theta+\varphi),
			\end{aligned}
			\]
			其中 $\varphi$ 由 $\cos\varphi=\frac{12}{\sqrt{369}},\sin\varphi=\frac{15}{\sqrt{369}}$ 决定。由于 $20>3\sqrt{41}$，故 $|15\cos\theta+12\sin\theta-20| = 20-M$，所以 $d=\frac{20-M}{5}$。当 $M=-3\sqrt{41}$ 时，$d_{\max}=\frac{20+3\sqrt{41}}{5}$。
			\par\textbf{理由：} 三角函数有界性及绝对值处理。
	\end{description}
	\textbf{\textcolor{CExample}{关键结论：}}
	\textbf{最大距离为 $\frac{20+3\sqrt{41}}{5}$。}

	\medskip
	\textbf{\textcolor{CExample}{例 3（综合应用）}}
	\textbf{\textcolor{CExample}{题目：}}
	已知椭圆 $C:\frac{x^2}{9}+\frac{y^2}{4}=1$ 和定点 $A(1,0)$，点 $P$ 在椭圆上运动，求 $|PA|$ 的最大值和最小值。
	\textbf{\textcolor{CExample}{解题步骤：}}
	\begin{description}[style=nextline, leftmargin=2.8em, labelsep=0.5em, itemsep=0.45em]
		\item[\textbf{第一步（参数设点）：}]
			设 $P(3\cos\theta,2\sin\theta)$。
			\par\textbf{理由：} 参数方程表示椭圆上点。
		\item[\textbf{第二步（表达距离平方）：}]
			\[
				\begin{aligned}
					|PA|^2 &= (3\cos\theta-1)^2+(2\sin\theta)^2 \\ &= 9\cos^2\theta-6\cos\theta+1+4\sin^2\theta \\ &= 5\cos^2\theta-6\cos\theta+5.
			\end{aligned}
			\]
			\par\textbf{理由：} 展开并利用 $\sin^2\theta=1-\cos^2\theta$ 化简。
		\item[\textbf{第三步（换元求最值）：}]
			令 $t=\cos\theta\in[-1,1]$，则 $|PA|^2=5t^2-6t+5$。对称轴 $t=\frac{3}{5}\in[-1,1]$。 最小值：当 $t=\frac{3}{5}$ 时，
			\[
				\begin{aligned}
					|PA|^2_{\min} &= 5\cdot\left(\frac{3}{5}\right)^2-6\cdot\frac{3}{5}+5 \\ &= \frac{16}{5}.
			\end{aligned}
			\]
			最大值：当 $t=-1$ 时，
			\[
				\begin{aligned}
					|PA|^2_{\max} &= 5\cdot(-1)^2-6\cdot(-1)+5 \\ &= 16.
			\end{aligned}
			\]
			因此 $|PA|_{\min} = \sqrt{\frac{16}{5}}$，即 $|PA|_{\min} = \frac{4\sqrt5}{5}$； $|PA|_{\max}=4$。
			\par\textbf{理由：} 二次函数在闭区间上的最值。
	\end{description}
	\textbf{\textcolor{CExample}{关键结论：}}
	\textbf{$|PA|_{\min}=\frac{4\sqrt5}{5}$，$|PA|_{\max}=4$。}
\end{examplebox}
